% Tokyo-To administrative hierarchy (schematic). pdfLaTeX-safe (Latin labels only).
\begin{figure}[htbp]
\centering
\begin{tikzpicture}[
  node distance=1.15cm,
  box/.style={rectangle, draw=black!70, rounded corners=2pt, align=center,
    minimum width=6.2cm, minimum height=1.05cm, font=\small},
  arrow/.style={-{Stealth[length=2.2mm]}, thick, black!75}
]
\node[box] (to) {Tokyo Metropolitan Government (\emph{Tokyo-To})};
\node[box, below=of to] (kyoku) {Kyoku --- bureau / department};
\node[box, below=of kyoku] (bu) {Bu --- division (when listed)};
\node[box, below=of bu] (ka) {Ka --- section};
\node[box, below=of ka] (kakari) {Kakari --- smallest unit (team / desk)};
\draw[arrow] (to) -- (kyoku);
\draw[arrow] (kyoku) -- (bu);
\draw[arrow] (bu) -- (ka);
\draw[arrow] (ka) -- (kakari);
\node[below=0.35cm of kakari, font=\footnotesize, text width=11.5cm, align=center]
  {Headings follow the printed yearbooks; not every branch lists all four levels.
  The analysis forward-fills these fields so each worker maps to a \emph{ka} and \emph{kakari} for fixed effects.};
\end{tikzpicture}
\caption{Administrative hierarchy in Tokyo metropolitan government (\emph{Tokyo-To}) personnel directories. The diagram is schematic; some listings omit a \emph{bu} line or combine labels.}
\label{fig:tokyo-hierarchy}
\end{figure}
